\chapter{Benchmarks}

\section{Internal benchmark: D-Stress workload}
256 independent neurons, 128 inputs each (16 tiles of \code{P\_IN}=8),
one shared activation vector -- the project's own existing, reused
benchmark, run against the final, frozen RTL via Verilator.

\begin{tabularx}{\textwidth}{L{2cm}L{2.4cm}L{2.4cm}Y}
\toprule
\rowh \thd{Config} & \thd{Cycles} & \thd{Correctness} & \thd{Latency @ 64\,MHz} \\
\midrule
N=2 & 49\,788 & 256/256 \PASS & 777.9\,\textmu s \\
\rowa N=4 & 49\,771 & 256/256 \PASS & 777.7\,\textmu s \\
\bottomrule
\end{tabularx}

Cycle counts are \textbf{simulation-verified} (Verilator, bit-exact)
and unchanged before/after the ERR-0025 read-timing fix. Latency at
64\,MHz is a \textbf{derived} figure (cycles~/~clock), not an
independently re-simulated wall-clock measurement, and is not a
measured-hardware number (no fabricated board exists).

N=4 barely improves over N=2 on this specific workload
($49\,771$ vs.\ $49\,788$ cycles) because the workload is dominated by
shared-SDRAM and shared-activation-fill traffic, not by raw MAC
throughput -- consistent with this project's own prior STEP17/18
memory-bound-vs-compute-bound findings, not a new anomaly.

\section{Host-interface pacing benchmark}
The board-level SPI smoke test (\S\ref{sec:host}) exercises single-job,
back-to-back, and gap-swept (100\,ns/5\,000\,ns/50\,000\,ns/$\sim$85\,\textmu s)
dispatch patterns -- all 11/11 bit-exact -- establishing that realistic
host pacing does not, itself, change correctness (only the STEP19-era
registered-read bug did, and that is fixed).

\section{Software reference comparison}
\textbf{No physical ESP32 (or other embedded MCU) hardware is
available in this environment.} A real ESP32 benchmark was
\emph{not} performed, and no ESP32 number is estimated or invented
here -- this is a disclosed OPEN item, not a claimed result.

As an honest, clearly-labeled illustrative data point only, the
identical D-Stress arithmetic (256 neurons $\times$ 128 INT8 MACs,
ReLU + INT8 saturate) was compiled (\code{cc~-O2}) and run on the
\emph{development machine itself} (Apple M4, arm64 -- \textbf{not}
an embedded target):

\begin{tabularx}{\textwidth}{L{4.4cm}Y}
\toprule
\rowh \thd{Metric} & \thd{Value} \\
\midrule
Platform & Apple M4 (arm64), \code{cc -O2}, single-threaded scalar C \\
\rowa Latency, full 256-neuron batch & 1.28\,\textmu s \\
Throughput & $2.56\times10^{10}$ MAC/s \\
\bottomrule
\end{tabularx}

On this specific, tiny workload, a modern, auto-vectorizing,
multi-GHz superscalar CPU outperforms the current FPGA design's
\emph{simulated} cycle count in raw latency. This is reported
honestly rather than omitted: FPGA-Neural V2's real value proposition
is a deterministic, low-power, standalone SDRAM-attached accelerator
for hosts that do \emph{not} have this class of CPU (e.g.\ an
ESP32-class microcontroller) -- not a claim of outperforming a
desktop-class processor on this workload size. A real, meaningful
software baseline for THAT comparison requires the actual embedded
target hardware, which remains unavailable this revision.
